\section{Introduction}
\textsc{Phi} is a constructed language. A constructed language is one whose phonology, grammar, and vocabulary, instead of having developed naturally, are consciously devised\footnote{\hspace{1.1em}\url{https://en.wikipedia.org/wiki/Constructed\_language}}. The sound of the name \emph{Phi} embodies the simplicity and softness of the language itself. \emph{Phi~($\varPhi$)} is also beauty in mathematics. It is

\begin{quote}
  \emph{the golden ratio, an irrational number with a value of approximately 1.618033988 which expresses the relationship that the sum of two quantities is to the larger quantity as the larger is to the smaller.} -- Wiktionary \footnote{\url{https://en.wiktionary.org/wiki/\%CF\%86}}
\end{quote}

It is thought that the value of \emph{phi} is one of the fundamental characteristics of nature; in the language Phi, nature is fundamental to its being. Listening to spoken Phi one is sometimes reminded of the sounds of the natural world. There is also a musicality and rhythm to sentences in the language.
